% ===================================================================
%  Αρχείο: 37086_SOLUTION.tex (Fragment)
%  Αυτόματη μετατροπή μέσω doc2latex.py
% ===================================================================

ΘΕΜΑ 4

Θεωρούμε παραλληλόγραμμο ΑΒΓΔ και τις προβολές Α΄, Β΄, Γ΄, Δ΄ και Ο΄ των κορυφών του Α, Β, Γ, Δ και του κέντρου του Ο αντίστοιχα, σε μία ευθεία ε.

\begin{alist}
\item Αν η ευθεία (ε) αφήνει τις κορυφές του παραλληλογράμμου στο ίδιο ημιεπίπεδο (όπως στο σχήμα που ακολουθεί) και είναι ΑΑ΄= 3, ΒΒ΄= 2, ΓΓ΄= 5, τότε:

.}
\end{alist}
\begin{alist}
\item
  Να αποδείξετε ότι η απόσταση ΟΟ' του κέντρου Ο του παραλληλογράμμου από την (ε) είναι ίση με 4. \monades{8}
\item
  Να βρείτε την απόσταση ΔΔ΄. \monades{9}


% TODO: Μετατροπή σχήματος σε TikZ/PGFPlots
\begin{figure}[htbp]
\centering
\includegraphics[width=0.8\linewidth]{../extracted/37086_SOLUTION/image1.png}
% \caption{Σχήμα 1}
\end{figure}


\item Αν η ευθεία (ε) διέρχεται από το κέντρο του παραλληλογράμμου και είναι παράλληλη προς δύο απέναντι πλευρές του, τι παρατηρείτε για τις αποστάσεις ΑΑ΄, ΒΒ΄, ΓΓ΄, ΔΔ΄.

Να αιτιολογήσετε την απάντησή σας. \monades{8}
\end{alist}
